% gerado por scripts/build_paper_assets.py -- nao editar a mao
\begin{table}[htbp]
\centering
\small
\setlength{\tabcolsep}{5pt}
\caption{Effect of pooling the eight cities into a single model: difference in RelMAE
between the global and the per-city form of the same estimator.}
\label{tab:global}
\begin{tabular}{lccccc}
\toprule
City & CatBoost & LightGBM & Random Forest & Ridge & XGBoost \\
\midrule
Sao Paulo & +0.03 & -0.03 & -0.05 & -0.25 & -0.09 \\
Rio de Janeiro & +0.39 & +0.46 & -1.51 & -0.26 & -1.94 \\
Belo Horizonte & -0.02 & -0.88$^{\ast}$ & -0.27 & -0.06 & -0.25 \\
Brasilia & +0.02 & -0.12 & -0.06 & -0.05 & -0.07 \\
Fortaleza & -0.14 & -0.62 & -0.97$^{\ast}$ & -0.22$^{\ast}$ & -0.85$^{\ast}$ \\
Recife & -0.24 & -0.49$^{\ast}$ & +0.20 & -0.16 & +0.25 \\
Manaus & -0.05 & +0.81 & -0.25$^{\ast}$ & +0.76$^{\ast}$ & -0.09 \\
Salvador & +0.04 & +0.04 & +0.08 & -0.09 & +0.01 \\
\bottomrule
\end{tabular}

\vspace{0.5em}
\begin{minipage}{\textwidth}
\footnotesize
Each cell is the RelMAE of the global model minus the RelMAE of the per-city model fitted
with the same estimator, so a negative value means that pooling the eight series helped.
Global models were fitted on all eight cities jointly with per-series standardisation and
the city identity entered as an indicator. Comparing a global model of one estimator with
a per-city model of another would confound pooling with the choice of learner, so every
comparison is within an estimator. An asterisk marks the pairs that satisfy the same
pre-registered decision rule applied to the baseline comparison, which is defined in the
Methods. This table is reproduced in both the main text and the supplementary material, so
it does not cross-reference a table number that exists in only one of them. Pooling helped in 6 of the
40 city-estimator pairs and harmed in 1; in the remaining
33 the test did not distinguish the two forms. No global model
beat the seasonal naive baseline in any city.
\end{minipage}
\end{table}
